\documentclass{szb-solution}

\area{Sorozatok}
\title{Numerikus sorozatok II}
\subject{Matematika G1}
\subjectCode{BMETE94BG01}
\date{Utoljára frissítve: \today}
\docno{5}

\begin{document}
\maketitle

\subsection{Konvergencia a definíció alapján}

\begin{enumerate}[a)]
  \item Legyen
        $$
          a_n=\frac{2n+5}{n-1}
          \text.
        $$
        A számláló és a nevező legmagasabb fokú tagjainak együtthatóiból
        sejthető, hogy $a_n\to2$. Ezt a definícióval is igazoljuk:
        $$
          |a_n-2|
          =\left|\frac{2n+5-2(n-1)}{n-1}\right|
          =\frac{7}{n-1}
          \qquad(n>1)
          \text.
        $$
        Adott $\varepsilon>0$ esetén a
        $$
          \frac{7}{n-1}<\varepsilon
        $$
        egyenlőtlenség teljesül, ha $n>1+7/\varepsilon$. Választhatjuk tehát
        $$
          N(\varepsilon)=\left\lceil 1+\frac7\varepsilon\right\rceil
          \text.
        $$
        Az $\varepsilon=10^{-6}$ értékhez például $N=7\,000\,001$
        megfelelő. Ekkor $n>N$ esetén $|a_n-2|<10^{-6}$, vagyis
        $$
          \boxed{a_n\longrightarrow2}
          \text.
        $$

  \item Legyen
        $$
          b_n=\frac{2n+3\sqrt n}{3n+1}
          \text.
        $$
        A várható határérték $2/3$, és
        \begin{align*}
          \left|b_n-\frac23\right|
          &=\left|\frac{3(2n+3\sqrt n)-2(3n+1)}{3(3n+1)}\right|\\
          &=\frac{9\sqrt n-2}{9n+3}
          <\frac{9\sqrt n}{9n}
          =\frac1{\sqrt n}
          \qquad(n\geq1)
          \text.
        \end{align*}
        Ezért elég $n>1/\varepsilon^2$-et megkövetelni. Az
        $\varepsilon=10^{-3}$ értékhez például $N=1\,000\,000$ választható.
        (Ez egyszerű, biztonságos választás; nem a legkisebb lehetséges
        küszöb.) Így
        $$
          \boxed{b_n\longrightarrow\frac23}
          \text.
        $$
\end{enumerate}

\subsection{Torlódási pontok}

\begin{enumerate}[a)]
  \item Írjuk a sorozatot két tényező szorzataként:
        $$
          a_n=
          \underbrace{\frac{(\sqrt{n^2+1}+n)^2}{\sqrt[3]{n^6+1}}}_{u_n}
          \cos\frac{n\pi}{3}
          \text.
        $$
        Az első tényezőt $n^2$-tel normálva kapjuk, hogy
        $$
          u_n=
          \frac{(\sqrt{1+1/n^2}+1)^2}{\sqrt[3]{1+1/n^6}}
          \longrightarrow4
          \text.
        $$
        A második tényező $6$-periodikus. Az $n$ maradéka szerint:
        $$
          \cos\frac{n\pi}{3}\in
          \left\{1,\frac12,-\frac12,-1\right\}
          \text,
        $$
        és mind a négy érték végtelen sokszor előfordul. Az egyes
        maradékosztályokhoz tartozó részsorozatok határértékei tehát
        $4$, $2$, $-2$ és $-4$. Más torlódási pont nincs, mert minden index
        a hat maradékosztály egyikébe tartozik. Így
        $$
          \boxed{T(a_n)=\{-4,-2,2,4\}}
          \text.
        $$

  \item Egyszerűsítés után
        $$
          b_n=\frac{(-1)^n+2}{3}\sin\frac{2n\pi}{3}
          \text.
        $$
        Az első tényező páros $n$ esetén $1$, páratlan $n$ esetén $1/3$;
        a szinuszos tényező pedig $3$-periodikusan a
        $\sqrt3/2$, $-\sqrt3/2$, $0$ értékeket veszi fel. A teljes sorozat
        ezért $6$-periodikus. A hat maradékosztály értékei rendre
        $$
          \frac{\sqrt3}{6},\quad -\frac{\sqrt3}{2},\quad 0,
          \quad \frac{\sqrt3}{2},\quad -\frac{\sqrt3}{6},\quad 0
          \text.
        $$
        Mivel egy periodikus sorozat minden, a periódusban előforduló értéke
        torlódási pont,
        $$
          \boxed{
            T(b_n)=\left\{-\frac{\sqrt3}{2},-\frac{\sqrt3}{6},0,
            \frac{\sqrt3}{6},\frac{\sqrt3}{2}\right\}
          }
          \text.
        $$
\end{enumerate}

\subsection{Monotonitás és korlátosság}

\begin{enumerate}[a)]
  \item Legyen
        $$
          a_n=\frac{2n^2+1}{n^2-n+1}
          \text.
        $$
        A szomszédos tagok különbsége
        \begin{align*}
          a_{n+1}-a_n
          &=\frac{2n^2+4n+3}{n^2+n+1}
            -\frac{2n^2+1}{n^2-n+1}\\
          &=\frac{2(1-n^2)}{(n^2+n+1)(n^2-n+1)}
          \text.
        \end{align*}
        A nevező pozitív. $n=1$ esetén a különbség nulla, $n\geq2$ esetén
        negatív. Tehát $a_1=a_2=3$, majd a sorozat szigorúan csökken.
        Továbbá
        $$
          a_n-2=\frac{2n-1}{n^2-n+1}>0,
          \qquad
          \lim_{n\to\infty}a_n=2
          \text.
        $$
        Következésképpen
        $$
          \boxed{2<a_n\leq3}
          \text,
        $$
        vagyis a sorozat korlátos, és $a_n\to2$.

  \item Legyen $b_n=n^5/n!$. A hányadosvizsgálat:
        $$
          \frac{b_{n+1}}{b_n}
          =\frac{(n+1)^4}{n^5}
          \text.
        $$
        Az első három hányados $16$, $81/32$, illetve $256/243$, tehát a
        sorozat $b_1$-től $b_4$-ig nő. Ha $n\geq4$, akkor
        $$
          \frac{(n+1)^4}{n^5}
          =\frac1n\left(1+\frac1n\right)^4
          \leq\frac14\left(\frac54\right)^4
          =\frac{625}{1024}<1
          \text,
        $$
        ezért $b_4$ után szigorúan csökken. A teljes sorozat tehát nem
        monoton, de pozitív és a maximuma
        $$
          b_4=\frac{4^5}{4!}=\frac{128}{3}
          \text.
        $$
        A fenti becslésből $b_{n+1}\leq(625/1024)b_n$ következik $n\geq4$
        esetén, így a farokrész egy nullához tartó mértani sorozattal
        becsülhető. Ezért
        $$
          \boxed{0<b_n\leq\frac{128}{3},\qquad b_n\longrightarrow0}
          \text.
        $$
\end{enumerate}

\subsection{Lineáris rekurzió}

A rekurzió ismételt alkalmazásával
$$
  a_n=1+\frac12+\frac1{2^2}+\dots+\frac1{2^{n-1}}
  \text.
$$
Ez egy véges mértani összeg, ezért
$$
  a_n=\frac{1-(1/2)^n}{1-1/2}
  =2-\frac1{2^{n-1}}
  \text.
$$
Az explicit alakból azonnal látszik, hogy a sorozat szigorúan növekvő,
$1\leq a_n<2$, és
$$
  \boxed{a_n\longrightarrow2}
  \text.
$$

\subsection{Gyökös rekurzió}

Legyen $a_1=1$ és $a_n=\sqrt{1+a_{n-1}}$. Jelölje
$$
  \varphi=\frac{1+\sqrt5}{2}
  \text,
$$
ekkor $\varphi^2=1+\varphi$.

Először belátjuk indukcióval, hogy $1\leq a_n<\varphi$. Ez $n=1$-re igaz.
Ha $a_n<\varphi$, akkor a gyökfüggvény szigorú monotonitása miatt
$$
  a_{n+1}=\sqrt{1+a_n}<\sqrt{1+\varphi}=\varphi
  \text.
$$
Az alsó korlát is nyilvánvaló, így az indukció teljes.

A növekedést szintén indukcióval igazoljuk. $a_2=\sqrt2>a_1$. Ha
$a_n>a_{n-1}$, akkor
$$
  a_{n+1}=\sqrt{1+a_n}>\sqrt{1+a_{n-1}}=a_n
  \text.
$$
Tehát $(a_n)$ növekvő és felülről korlátos, ezért konvergens. Ha a
határértéke $A$, akkor a gyökfüggvény folytonossága miatt
$$
  A=\sqrt{1+A},
  \qquad A^2-A-1=0
  \text.
$$
Mivel $A\geq1$, csak a pozitív gyök jöhet szóba:
$$
  \boxed{a_n\longrightarrow\frac{1+\sqrt5}{2}}
  \text.
$$

\subsection{Komplex sorozatok}

Komplex sorozat akkor és csak akkor konvergens, ha a valós és a képzetes
része is konvergens; ekvivalensen a különbség abszolút értéke nullához tart.

\begin{enumerate}[a)]
  \item $n^2$-tel osztva
        $$
          \frac{n^2-i(n^2-1)}{n^2-i}
          =\frac{1-i+i/n^2}{1-i/n^2}
          \longrightarrow\boxed{1-i}
          \text.
        $$

  \item A háromszög-egyenlőtlenségből
        $$
          \left|\frac{i^n}{3^n+i^n}\right|
          \leq\frac{|i^n|}{3^n-|i^n|}
          =\frac1{3^n-1}\longrightarrow0
          \text.
        $$
        Ezért
        $$
          \boxed{\frac{i^n}{3^n+i^n}\longrightarrow0}
          \text.
        $$

  \item
        $$
          |(1-i)^n|=|1-i|^n=(\sqrt2)^n\longrightarrow\infty
          \text.
        $$
        Egy konvergens komplex sorozat szükségképpen korlátos, ez a sorozat
        viszont nem az. Tehát
        $$
          \boxed{((1-i)^n)\text{ divergens}}
          \text.
        $$
\end{enumerate}

\subsection{Egy hatványegyenlőtlenség}

Az állítás $n\geq3$ esetén ekvivalens a
$$
  n\geq\left(1+\frac1n\right)^n
  \text.
$$
egyenlőtlenséggel. A binomiális tétel szerint
\begin{align*}
  \left(1+\frac1n\right)^n
  &=\sum_{k=0}^n\binom nk\frac1{n^k}
  =\sum_{k=0}^n
    \frac{n(n-1)\cdots(n-k+1)}{k!\,n^k}\\
  &\leq\sum_{k=0}^n\frac1{k!}
  <3
  \leq n
  \text.
\end{align*}
Itt az utolsó szigorú becsléshez felhasználtuk, hogy $k!\geq2^{k-1}$
$k\geq2$ esetén, és legalább egy helyen szigorú az egyenlőtlenség. Ezzel
$$
  \boxed{n^{n+1}\geq(n+1)^n\qquad(n\geq3)}
  \text.
$$

\subsection{Binomiális együtthatók hányadosa}

A sorozat $n\geq3$ esetén értelmezett, és
$$
  a_n=\frac{\binom n2}{\binom n3}
  =\frac{n!}{(n-2)!2!}\frac{(n-3)!3!}{n!}
  =\frac3{n-2}
  \text.
$$
Innen
$$
  a_{n+1}-a_n
  =\frac3{n-1}-\frac3{n-2}
  =-\frac3{(n-1)(n-2)}<0
  \text,
$$
tehát a sorozat szigorúan csökken. Továbbá
$$
  \boxed{0<a_n\leq3,\qquad a_n\longrightarrow0}
  \text.
$$

\end{document}
